% !TEX encoding = utf8
\chapter{Rysunki}
Wszystkie elementy dokumentu opracowanego w~systemie \LaTeX\ powinny wyglądać jednolicie. Do wykonywania rysunków korzystamy więc z mechanizmów oferowanych przez dodatkowe pakiety \LaTeX a, nie dołączamy rysunków wykonanych jakościowo różnych, np. wykonanych w~programie Word. Nie używamy rysunków zapisanych w~plikach bitmapowych, lecz w~plikach wektorowych (\verb+pdf+, ew. \verb+ps+ lub \verb+eps+). Jedynym wyjątkiem są zdjęcia.

\section{Schematy blokowe}
Do opracowania schematów blokowych najlepiej wykorzystać język opisu rysunków \verb+TikZ/PGF+ \cite{litTantau2015}. Przy wykonywaniu prostych rysunków po prostu opisujemy je za pomocą poleceń dodających kolejne elementy, tzn. prostokąty, okręgi, linie. Na przykład, ciąg poleceń:
\begin{lstlisting}[style=customlatex,frame=single] 
\begin{figure}[b]
\centering
\begin{tikzpicture}
\draw [red, thick] (0,0) rectangle (4,3);
\draw [orange, thick,dashed] (1,1) rectangle (3,2);
\draw [->,thick] (4,1.5) -- (7,1.5);
\draw [blue, thick] (8,1.5) circle [radius=1];
\node at (2,1.5) {Tekst};
\node at (8,1.5) {$y(x)=x^2$};
\end{tikzpicture}
\end{figure}
\end{lstlisting}
pozwala narysować figury geometryczne przedstawione na rys. \ref{r_tikz_przyklad}. Zwróćmy uwagę, że napis oraz wzór są złożone aktualnie wykorzystywaną czcionką, jej wielkość jest taka sama jak w~całym dokumencie.

\begin{figure}[b]
\centering
\begin{tikzpicture}
\draw [red, thick] (0,0) rectangle (4,3);
\draw [orange, thick,dashed] (1,1) rectangle (3,2);
\draw [->,thick] (4,1.5) -- (7,1.5);
\draw [blue, thick] (8,1.5) circle [radius=1];;
\node at (2,1.5) {Tekst};
\node at (8,1.5) {$y(x)=x^2$};
\end{tikzpicture}

\caption{Przykładowy rysunek wykonany w~języku \texttt{TikZ/PGF}}
\label{r_tikz_przyklad}
\end{figure}

Przy większych rysunkach można wykorzystać programy umożliwiające ich przygotowanie przy wykorzystaniu środowiska graficznego, np. \verb+TikzEdt+ (\url{http://www.tikzedt.org/}), \verb+TpX+ (\url{http://tpx.sourceforge.net/}), \verb+ktikz+ (\url{https://www.linux-apps.com/p/1126914/}), \verb+GraTeX+ (\url{https://sourceforge.net/projects/gratex/}).

Można również wykorzystać starsze, klasyczne pakiety \verb+picture+, \verb+epic+, \verb+eepic+. Również w~tych przypadkach można ,,ręcznie'' opisywać poszczególne elementy graficzne lub skorzystać ze środowiska graficznego, np. \verb+LaTeXPiX+ (\url{http://latexpix.software.informer.com/}), które znacznie przyspiesza pracę. Inne narzędzia, umożliwiające opracowanie rysunków wysokiej jakości, to \verb+METAPOST+ oraz \verb+PSTricks+.

\section{Kolory}
Przy umieszczaniu kilku wykresów na tym samym rysunku, np. wyników eksperymentów dla różnych wartości parametrów, należy zastosować kolory różniące się od siebie w~znacznym stopniu, nie można stosować kolorów podobnych, np. kilku odcieni tego samego koloru. Warto zastosować standardową paletę kolorów o~nazwie \texttt{lines}, użytą w~pakiecie MATLAB, która została przedstawiona w~tab.~\ref{r_zestaw_kolorow}. Definicja tych kolorów w~systemie \LaTeX\ jest następująca:
\begin{lstlisting}[style=customlatex,frame=single]
\definecolor{niebieski}{rgb}{0,0.447,0.741}
\definecolor{czerwony}{rgb}{0.85,0.325,0.098}	
\definecolor{zolty}{rgb}{0.929,0.694,0.125}		
\definecolor{fioletowy}{rgb}{0.494,0.184,0.556}	
\definecolor{zielony}{rgb}{0.466,0.674,0.188}
\definecolor{jasnoniebieski}{rgb}{0.301,0.745,0.933}
\definecolor{ciemnioczerwony}{rgb}{0.635,0.078,0.184}
\end{lstlisting}

\definecolor{niebieski}{rgb}{0,0.447,0.741}
\definecolor{czerwony}{rgb}{0.85,0.325,0.098}	
\definecolor{zolty}{rgb}{0.929,0.694,0.125}		
\definecolor{fioletowy}{rgb}{0.494,0.184,0.556}	
\definecolor{zielony}{rgb}{0.466,0.674,0.188}
\definecolor{jasnoniebieski}{rgb}{0.301,0.745,0.933}
\definecolor{ciemnioczerwony}{rgb}{0.635,0.078,0.184}

\begin{table}[h]
\caption{Paleta 7 kolorów}
\label{r_zestaw_kolorow}
	\centering
	\begin{small}
		\begin{tabular}{|p{3cm}|p{4cm}|}
			\hline
			%\multicolumn{1}{|c|}{Nazwa koloru\rule{0pt}{3.5mm}} & \multicolumn{1}{|c|}{} \\ \hline
			niebieski & \cellcolor{niebieski}\\
			czerwony & \cellcolor{czerwony}\\
			zolty & \cellcolor{zolty}\\
			fioletowy & \cellcolor{fioletowy}\\
			zielony & \cellcolor{zielony}\\
			jasnoniebieski & \cellcolor{jasnoniebieski}\\
			ciemnioczerwony & \cellcolor{ciemnioczerwony}\\
			\hline
		\end{tabular}
	\end{small}
\end{table}

\section{Funkcje statyczne}
Do wykonywania wykresów prezentujących wyniki symulacji i~eksperymentów stosuje się pakiet \verb+PGFPLOTS+ \cite{litFeuersanger2016}. Załóżmy, że w~katalogu \verb+rysunki/dane_stat+ znajduje się plik \verb+dane_fx.txt+ zawierający w~pierwszej kolumnie wartości argumentu $x$, natomiast w~drugiej kolumnie wartości funkcji $f(x)$:
\begin{lstlisting}[style=customlatex,frame=single]
  -10.0000 -808.7350
   -9.0000 -696.4791
   -8.0000 -539.7850
   -7.0000 -386.1268
   -6.0000 -291.5881
   -5.0000 -267.6436
   -4.0000 -268.9551
   -3.0000 -233.3995
   -2.0000 -137.5303
   -1.0000  -16.4788
         0   70.0000
    1.0000   94.1212
    2.0000   87.2697
    3.0000  112.8005
    4.0000  209.4449
    5.0000  357.3564
    6.0000  498.0119
    7.0000  589.6732
    8.0000  647.4150
    9.0000  730.9209
   10.0000  891.2650
\end{lstlisting}
oraz podobny plik \verb+dane_gx.txt+, definiujący funkcję $g(x)$. Aby narysować te funkcje stosuje się polecenia:
\begin{lstlisting}[style=customlatex,frame=single]
\newcommand{\szer}{8cm}
\newcommand{\wys}{7cm}

\begin{figure}[t]
\centering
\begin{tikzpicture}
\begin{axis}[
	width=\szer,
	height=\wys,
	xmin=-10,xmax=10,ymin=-1000,ymax=1000,
	xlabel={$x$},
	ylabel={$f(x), \ g(x)$},
	xtick={-10,-5,0,5,10},
	ytick={-1000,-500,0,500,1000},
	legend pos=south east,
]

\definecolor{niebieski}{rgb}{0,0.447,0.741}
\definecolor{czerwony}{rgb}{0.85,0.325,0.098}	

\addplot[niebieski,thick]
	file {rysunki/dane_stat/dane_fx.txt};
\addplot[czerwony,thick,densely dashed]
	file {rysunki/dane_stat/dane_gx.txt};
\legend{$f(x)$,$g(x)$}
\end{axis}
\end{tikzpicture}
\caption{Przykładowy rysunek funkcji $f(x)$ i~$g(x)$ wykonany 
w~języku \texttt{PGFPLOTS}}
\label{r_pgfplots_funkcje}
\end{figure}
\end{lstlisting}
Otrzymany rezultat przedstawiono na rys.~\ref{r_pgfplots_funkcje}.

Dla wykorzystanej palety kolorów ładnie wyglądają rysunki o~pogrubionych liniach (styl \verb+thick+).

Istnieje możliwość ustawienia wielkości czcionek liczb umieszczonych: na osiach (\verb+tick label style+), w~oznaczeniach osi (\verb+label style+), w~legendzie (\verb+legend style+) oraz w~tytule rysunku (\verb+title style+). Przykładowa konfiguracja zmieniająca wielkość czcionek jest następująca:
\begin{lstlisting}[style=customlatex,frame=single]
\pgfplotsset{
	tick label style={font=\tiny},
	label style={font=\footnotesize},
	legend style={font=\footnotesize},
	title style={font=\footnotesize},
	/pgf/number format/.cd, use comma, 1000 sep={}
}
\end{lstlisting}
Dodatkowo, sekwencja \verb+/pgf/number format/.cd, use comma, 1000 sep={}+ ustawia przecinek jako separator dziesiętny (zamiast domyślnej kropki) oraz kasuje separator, oddzielający tysiące od setek (domyślenie jest to przecinek).

\newcommand{\szer}{8cm}
\newcommand{\wys}{7cm}

\begin{figure}[tb]
\centering
\begin{tikzpicture}
\begin{axis}[
	width=\szer,
	height=\wys,
	xmin=-10,xmax=10,ymin=-1000,ymax=1000,
	xlabel={$x$},
	ylabel={$f(x), \ g(x)$},
	xtick={-10,-5,0,5,10},
	ytick={-1000,-500,0,500,1000},
	legend pos=south east,
]

\definecolor{niebieski}{rgb}{0,0.447,0.741}
\definecolor{czerwony}{rgb}{0.85,0.325,0.098}	

\addplot[niebieski,thick]
	file {rysunki/dane_stat/dane_fx.txt};
\addplot[czerwony,thick,densely dashed]
	file {rysunki/dane_stat/dane_gx.txt};
\legend{$f(x)$,$g(x)$}
\end{axis}
\end{tikzpicture}
\caption{Przykładowy rysunek funkcji $f(x)$ i~$g(x)$ wykonany w~języku \texttt{PGFPLOTS} (rysunek jest generowany przy każdym przetworzenia pliku źródłowego}
\label{r_pgfplots_funkcje}
\end{figure}

Opisany sposób implementacji rysunków jest poprawny, ale ma poważną wadę, ponieważ \LaTeX{}   potrzebuje dość dużo czasu na ich przetworzenie. Okazuje się to dużym mankamentem szczególnie wówczas, gdy w~dokumencie znajduje się dużo skomplikowanych rysunków. Skutecznym rozwiązaniem jest przygotowanie rysunków i~zapis ich do plików \verb+pdf+, a~następnie dołączenie ich do głównego dokumentu poleceniem \verb+\includegraphics+. W~katalogu \verb+rysunki/zapisz_pdf+ podano kilka przykładów takich rysunków. Plik źródłowy \verb+funkcje.tex+ ma postać:
\lstinputlisting[style=customlatex,frame=single]{rysunki/zapisz_pdf/funkcje.tex}
Zwróćmy uwagę, że dzięki zastosowaniu klasy \verb+standalone+, wygenerowany zostanie rysunek o~zadanych wymiarach (bez niepotrzebych marginesów), plik \verb+pdf+ nie będzie formatu A4. Następujące polecenie 
\begin{lstlisting}[style=customlatex,frame=single]
pdflatex funkcje.tex
\end{lstlisting}
zapisuje plik \verb+funkcje.pdf+. Wygenerowany rysunek dołącza się do dokumentu ciągiem instrukcji:
\begin{lstlisting}[style=customlatex,frame=single]
\begin{figure}[tb]
\centering
\includegraphics[scale=1]{pgfplots_pdf/funkcje}
\caption{Przykładowy rysunek funkcji $f(x)$ i~$g(x)$ wykonany
w~języku \texttt{PGFPLOTS}}
\label{r_pgfplots_funkcje}
\end{figure}
\end{lstlisting}
Otrzymany rezultat przedstawiono na rys.~\ref{r_pgfplots_funkcje_pdf}. Oczywiście, rysunki \ref{r_pgfplots_funkcje} i~\ref{r_pgfplots_funkcje_pdf} są bardzo podobne, inna jest tylko wielkość czcionek na osiach.

Przy dużych zbiorach danych \LaTeX{} zgłasza błąd pamięci. Należy wówczas zastosować system składu \hbox{Lua\LaTeX}.

Nie należy skalować rysunku przy wykorzystaniu argumentów polecenia \verb+\includegraphics+, gdyż zmieni to wielkość zastosowanej czcionki. Jeżeli zachodzi konieczność zmiany wielkości rysunku, należy zmodyfikować plik źródłowy generujący rysunek.

Jeżeli rysunek jest znacznie szerszy niż szerokość strony (oraz gdy nie możemy lub nie chcemy redukować szerokości rysunku), należy zastosować otoczenie \verb+sidewaysfigure+ z~pakietu \verb+rotating+, które działa analogicznie jak otoczenie \verb+sidewaystable+.

W~dokumentacji pakietu \verb+PGFPLOTS+ omówiono sposób przygotowania rysunków trójwymiarowych \cite{litFeuersanger2016}.

\begin{figure}[ptb]
\centering
\includegraphics[scale=1]{rysunki/zapisz_pdf/funkcje}
\caption{Przykładowy rysunek funkcji $f(x)$ i~$g(x)$ wykonany
w~języku \texttt{PGFPLOTS} i~zapisany w~pliku \texttt{funkcje.pdf}}
\label{r_pgfplots_funkcje_pdf}
\end{figure}

\section{Wyniki symulacji i~eksperymentów}
\subsection{Procesy jednowymiarowe}
Załóżmy, że w~katalogu \verb+rysunki/symulacje11+ znajduje się plik \verb+yzad.txt+ zawierający w~pierwszej kolumnie pomiary czasu $t$ (w~sekundach), natomiast w~drugiej kolumnie próbki sygnału wartości zadanej $y^{\mathrm{zad}}$. Wykonano symulacje algorytmu regulacji GPC przy pięciu różnych wartościach parametru $\lambda$: $\num{0,1}$, $\num{0,2}$, $\num{0,5}$, $\num{1}$ i~$\num{2}$. Przebiegi sygnału sterującego $u$ zapisano w~plikach \verb+u_lambda_0_1.txt+, \verb+u_lambda_0_2.txt+, \verb+u_lambda_0_5.txt+, \verb+u_lambda_1.txt+ i~\verb+u_lambda_2.txt+, natomiast przebiegi sygnału wyjściowego procesu $y$ zapisano w~plikach \verb+y_lambda_0_1.txt+, \verb+y_lambda_0_2.txt+, \verb+y_lambda_0_5.txt+, \verb+y_lambda_1.txt+ i~\verb+y_lambda_2.txt+. Przygotowano plik \verb+symulacje11.tex+, umożliwiający zapisanie rysunku do pliku \verb+symulacje11.pdf+. Znajduje się on w~katalogu \verb+rysunki/zapisz_pdf+ i~ma następującą postać:
\lstinputlisting[style=customlatex,frame=single]{rysunki/zapisz_pdf/symulacje11.tex}
Na początku powyższego pliku zostaje wczytany i~skonfigurowany pakiet \verb+siunitx+. Między innymi, ustawiamy przecinek jako separator dziesiętny. Liczby z~częścią ułamkową zapisujemy w~postaci \verb+\num{0.1}+. W~zależności od konfiguracji, otrzymamy $0.1$ lub $0{,}1$. Polecenie
\begin{lstlisting}[style=customlatex,frame=single]
pdflatex symulacje11.tex
\end{lstlisting}
zapisuje plik \verb+symulacje11.pdf+. Rezultat przedstawiono na rys.~\ref{r_pgfplots_symulacje11_pdf}. Rysunek dołącza się do dokumentu instrukcją \verb+\includegraphics+. Zwróćmy uwagę, że do narysowania wyników symulacji dla kolejnych wartości parametru $\lambda$ zastosowano różne kolory oraz różne style linii (linia ciągła, linia przerywana, itd.). Umożliwia to łatwe rozróżnienie dokumentów na czarno-białym wydruku. Przy wydruku kolorowym oraz dokumentach elektronicznych można zrezygnować ze stosowania różnych stylów linii, do ich rozróżnienia wystarczające są kolory, pod warunkiem jednak, że kolejne krzywe nie są położone bardzo blisko siebie.

Czasami, ze względu na ograniczoną objętość dokumentu, należy zmniejszyć rysunki. Aby zmniejszyć objętość można zastosować ułożenie poziome dwóch rysunków, prezentujących wyniki symulacji procesu. Przygotowano plik \verb+symulacje11_wersja2.tex+, umożliwiający zapisanie rysunku do pliku \verb+symulacje11_wersja2.pdf+. Znajduje się on w~katalogu \verb+rysunki/zapisz_pdf+ i~ma następującą postać:
\lstinputlisting[style=customlatex,frame=single]{rysunki/zapisz_pdf/symulacje11_wersja2.tex}
Rezultat przedstawiono na rys.~\ref{r_pgfplots_symulacje11_wersja2_pdf}.

\begin{figure}[H]
\centering
\includegraphics[scale=1]{rysunki/zapisz_pdf/symulacje11}
%\includegraphics[scale=1]{rysunki/zapisz_pdf/symulacje11}
\caption{Przykładowy rysunek wyników symulacji procesu jednowymiarowego wykonany
w~języku \texttt{PGFPLOTS} i~zapisany w~pliku \texttt{symulacje11.pdf}}
\label{r_pgfplots_symulacje11_pdf}
\end{figure}

\begin{figure}[H]
\centering
\includegraphics[scale=1]{rysunki/zapisz_pdf/symulacje11_wersja2}
\caption{Przykładowy rysunek wyników symulacji procesu jednowymiarowego wykonany
w~języku \texttt{PGFPLOTS} i~zapisany w~pliku \texttt{symulacje11\_wersja2.pdf}}
\label{r_pgfplots_symulacje11_wersja2_pdf}
\end{figure}

Warto podkreślić, że linie różne rodzaje linii, np. przerywane lub kropkowane, stosujemy wówczas, gdy poszczególne krzywe na wykresie położone są blisko siebie lub wręcz na siebie nachodzą. Jeżeli jednak poszczególne krzywe są od siebie oddalone, nie ma powodów aby nie stosować linii ciągłych, co poprawia czytelność rysunków. Na rys. \ref{r_pgfplots_symulacje11_linie_ciagle_pdf} i~\ref{r_pgfplots_symulacje11_wersja2_linie_ciagle_pdf} zamieszczono wersje rys. \ref{r_pgfplots_symulacje11_pdf} i~\ref{r_pgfplots_symulacje11_wersja2_pdf} wykonane przy użyciu linii ciągłych. Pliki służace do wygenerowania tych wykresów znajdują się w~katalogu \verb+rysunki/zapisz_pdf+

\begin{figure}[H]
\centering
\includegraphics[scale=1]{rysunki/zapisz_pdf/symulacje11_linie_ciagle}
\caption{Przykładowy rysunek wyników symulacji procesu jednowymiarowego wykonany
w~języku \texttt{PGFPLOTS} i~zapisany w~pliku \texttt{symulacje11\_linie\_ciagle.pdf}}
\label{r_pgfplots_symulacje11_linie_ciagle_pdf}
\end{figure}

\begin{figure}[H]
\centering
\includegraphics[scale=1]{rysunki/zapisz_pdf/symulacje11_wersja2_linie_ciagle}
\caption{Przykładowy rysunek wyników symulacji procesu jednowymiarowego wykonany
w~języku \texttt{PGFPLOTS} i~zapisany w~pliku \texttt{symulacje11\_wersja2\_linie\_ciagle.pdf}}
\label{r_pgfplots_symulacje11_wersja2_linie_ciagle_pdf}
\end{figure}

\subsection{Procesy wielowymiarowe}
Załóżmy, że w~katalogu \verb+rysunki/symulacje22+ znajdują się wyniki symulacji procesu o~dwóch wejściach i~dwóch wyjściach zapisane w~plikach: \verb+yzad1.txt+, \verb+yzad2.txt+, \verb+u1.txt+, \verb+u2.txt+, \verb+y1.txt+, \verb+y2.txt+. W~pierwszej kolumnie tych plików podano czas $t$ (w~sekundach), natomiast w~drugiej kolumnie wartość odpowiedniej zmiennej. Plik \verb+symulacje22.tex+, umożliwiający zapisanie rysunku do pliku \verb+symulacje22.pdf+ znajduje się w~katalogu \verb+rysunki/zapisz_pdf+ i~ma następującą postać:
\lstinputlisting[style=customlatex,frame=single]{rysunki/zapisz_pdf/symulacje22.tex}
Rezultat przedstawiono na rys.~\ref{r_pgfplots_symulacje22_pdf}.

Aby zmniejszyć objętość można nieco inaczej ułożyć 4 rysunki, prezentujące wyniki symulacji procesu dwuwymiarowego. Przygotowano plik \verb+symulacje22_wersja2.tex+, umożliwiający zapisanie rysunku do pliku \verb+symulacje22_wersja2.pdf+. Znajduje się on w~katalogu \verb+rysunki/zapisz_pdf+ i~ma następującą postać:
\lstinputlisting[style=customlatex,frame=single]{rysunki/zapisz_pdf/symulacje22_wersja2.tex}
Rezultat przedstawiono na rys.~\ref{r_pgfplots_symulacje22_wersja2_pdf}.

\begin{figure}[H]
\centering
\includegraphics[scale=1]{rysunki/zapisz_pdf/symulacje22}
\caption{Przykładowy rysunek wyników symulacji procesu dwuwymiarowego wykonany
w~języku \texttt{PGFPLOTS} i~zapisany w~pliku \texttt{symulacje22.pdf}}
\label{r_pgfplots_symulacje22_pdf}
\end{figure}

\begin{figure}[H]
\centering
\includegraphics[scale=1]{rysunki/zapisz_pdf/symulacje22_wersja2}
\caption{Przykładowy rysunek wyników symulacji procesu dwuwymiarowego wykonany
w~języku \texttt{PGFPLOTS} i~zapisany w~pliku \texttt{symulacje22\_wersja2.pdf}}
\label{r_pgfplots_symulacje22_wersja2_pdf}
\end{figure}

\section{Zapis wielu rysunków do jednego pliku \texttt{pdf}}
Aktywacja klasy \verb+standalone+ z~opcją \verb+tikz+ umożliwia zapis wielu rysunków do jednego pliku \texttt{pdf}. Rozważmy przykładowy plik \verb+funkcje2.tex+ z~katalogu \verb+rysunki/zapisz_pdf+:
\lstinputlisting[style=customlatex,frame=single]{rysunki/zapisz_pdf/funkcje2.tex}
Każdy rysunek wygenerowany przez otoczenie \verb+tikzpicture+ znajdzie się na osobnej stronie. Co ważne, jeżeli każdy rysunek będzie miał inne wymiary, dzięki klasie \verb+standalone+ każda strona jest przycinana do zawartości. Gotowe rysunki wstawiamy do dokumentu poleceniami (oczywiście w~otoczeniu \verb+figure+):
\begin{lstlisting}[style=customlatex,frame=single]
\includegraphics[page=1]{funkcje2.pdf}
\end{lstlisting}
oraz
\begin{lstlisting}[style=customlatex,frame=single]
\includegraphics[page=2]{funkcje2.pdf}
\end{lstlisting}
Opisany sposób zapisu rysunków pozwala na ograniczenie liczby plików oraz umożliwia grupowanie treści.

\section{Prosty eksport wykresów: skrypt \texttt{matlab2tikz}}
Oprócz opisanego sposobu przygotowywania rysunków z~wykorzystaniem pakietu \verb+PGFPLOTS+, można wykorzystać skrypt \texttt{matlab2tikz}, który generuje kod źródłowy (\verb+TikZ+) rysunku wykonanego w~programie MATLAB. Skrypt \texttt{matlab2tikz} dostępny jest pod adresem: \url{https://www.mathworks.com/matlabcentral/fileexchange/22022-matlab2tikz-matlab2tikz}. Rozważmy przykładowy plik \verb+wykresy_z_matlaba.m+ z~katalogu \verb+rysunki/matlab2tikz+:
\lstinputlisting[style=Matlab-editor]{./rysunki/matlab2tikz/wykresy_z_matlaba.m}
Wykres wstawiamy do dokumentu w~nstępujący sposób:
\begin{lstlisting}[style=customlatex,frame=single]
\begin{figure}[H]
\centering
	\input{./rysunki/matlab2tikz/wykres_tex.tex}
\caption{Przykładowy rysunek wykonany w~MATLABie 
i~wygenerowany skryptem \texttt{matlab2tikz}}
\label{r_matlab2tikz}
\end{figure}
\end{lstlisting}
Otrzymany rezultat przedstawiono na rys. \ref{r_matlab2tikz}. Oczywiście, czasami konieczna jest ingerencja i~drobna modyfikacja plików wygenerowanych przez skrypt \texttt{matlab2tikz}, np. w~celu zmiany wielkości czcionki lub całego rysunku.

Rysunek jest generowany na podstawie kodu źródłowego przy każdym przetworzeniu dokumentu przez system \LaTeX, a~więc przy większej liczbie rysunków celowe jest oddzielne przygotowanie poszczególnych rysunków w~plikach \verb+pdf+, a~następnie ich dołączenie do dokumentu przy użyciu polecenia \verb+\includegraphics+.

\begin{figure}[H]
\centering
	\input{./rysunki/matlab2tikz/wykres_tex.tex}
\caption{Przykładowy rysunek wykonany w~MATLABie i~wygenerowany skryptem \texttt{matlab2tikz}}
\label{r_matlab2tikz}
\end{figure}

Pomimo tego, że w~MATLAB-ie znakiem oddzielającym część całkowitą od części ułamkowej jest kropka, to w~sprawozdaniu zastąpiona ona została przecinkiem. Jest to wynikiem działania jednej z linijek znajdujących się w~preambule, które na stałe konfigurują pakiet PGFPlots tak, aby zawsze liczby formatował zgodnie z polskimi standardami.

Zwróćmy uwagę, że rysunek zapisany w~pliku źródłowym plik \verb+wykres_tex.tex+ jest generowany przy każdym przetworzeniu dokumentu przez system \LaTeX. Aby przyspieszyć działanie systemu, warto osobno wygenerować wszystkie rysunki, a~następnie wstawić do dokumentu gotowe pliki \verb+pdf+. Służy do tego plik \verb+wykres_tex_standalone.tex+, umożliwiający zapisanie rysunku do pliku \verb+wykres_tex_standalone.pdf+. Znajduje się on w~katalogu \verb+rysunki/matlab2tikz+ i~ma następującą postać:
\lstinputlisting[style=customlatex,frame=single]{rysunki/matlab2tikz/wykres_tex_standalone.tex}

W~celu efektywnej pracy ze skryptem \texttt{matlab2tikz} należy przenieść wszystkie potrzebne pliki do katalogu, w~którym MATLAB poszukuje wywoływanych funkcji. W~tym celu należy utworzyć katalog \verb|<MATLAB_ROOT>/toolbox/matlab2tikz| (gdzie \verb|<MATLAB_ROOT>| jest ścieżką do katalogu gdzie zainstalowany jest MATLAB), skopiować doń zawartość katalogu \verb|src|, a~następnie dodać powyższy katalog do listy ścieżek MATLAB-a (np. klikając zakładkę \verb|HOME| i~przycisk \verb|SetPath| w~MATLAB-ie).

\textbf{UWAGA:} MATLAB 2017b automatycznie dodaje do wykresów legendę (MATLAB 2017a tego nie robi). Dla zainteresowanych rozwiązaniem tego problemu, link\newline \url{https://github.com/matlab2tikz/matlab2tikz/issues/1006}

\subsection{MATLAB i~dobór kolorów}
Uwaga: nie stosujemy starej palety kolorów pakietu MATLAB, przedstawionej w~tab.~\ref{r_zestaw_kolorow_MATLAB_historia}.

\definecolor{staryniebieski}{rgb}{0,0,1}
\definecolor{staryczerwony}{rgb}{1,0,0}	
\definecolor{staryzolty}{rgb}{1,1,0}		
\definecolor{staryfioletowy}{rgb}{1,0,1}	
\definecolor{staryzielony}{rgb}{0,1,0}
\definecolor{staryjasnoniebieski}{rgb}{0,1,1}
\definecolor{czarny}{rgb}{0,0,0}

\begin{table}[H]
\caption{Stara paleta 7 kolorów pakietu MATLAB (nie używać)}
\label{r_zestaw_kolorow_MATLAB_historia}
	\centering
	\begin{small}
		\begin{tabular}{|p{3cm}|p{4cm}|}
			\hline
			%\multicolumn{1}{|c|}{Nazwa koloru\rule{0pt}{3.5mm}} & \multicolumn{1}{|c|}{} \\ \hline
			'r' & \cellcolor{staryczerwony}\\
			'g' & \cellcolor{staryzielony}\\
			'b' & \cellcolor{staryniebieski}\\
			'c' & \cellcolor{staryjasnoniebieski}\\
			'm' & \cellcolor{staryfioletowy}\\
			'y' & \cellcolor{staryzolty}\\
			'k' & \cellcolor{czarny}\\
			\hline
		\end{tabular}
	\end{small}
\end{table}

Aktualne wersje pakietu MATLAB automatycznie stosują kolejne kolory dla kolejnych poleceń \verb+plot+ lub \verb+stairs+ z~palety \texttt{lines}, przedstawione w~tab. \ref{r_zestaw_kolorow}. Nie musimy stosować żadnych identyfikatorów kolorów. Co ważne, nie stosujemy identyfikatorów kolorów postaci \verb+'r'+, \verb+'g'+, \verb+'b'+, \verb+'c'+, \verb+'m'+, \verb+'y'+, \verb+'k'+. Rozważmy plik \verb+kolory_matlab.m+ z~katalogu \verb+rysunki/matlab2tikz+:
\lstinputlisting[style=Matlab-editor]{./rysunki/matlab2tikz/kolory_matlab.m}
Powyższy skrypt generuje dwa rysunki, przy czym na pierwszym z~nich kolejne kolory są dobierane automatycznie, co osiągamy nie definiując żadnych identyfikatorów kolorów, na drugim rysunku używamy historycznej palety kolorów. Otrzymane rezultaty zamieszczono na rys. \ref{r_wykres_kolory1} i~rys. \ref{r_wykres_kolory2}. Kolory z~historycznej palety są zbyt jaskrawe, szczególnie zielony i~żółty. Reasumując, stosujemy automatyczny przydział kolorów (paleta \texttt{lines}).

\begin{figure}[H]
\centering
	\input{./rysunki/matlab2tikz/wykres_kolory1.tex}
\caption{Przykładowy rysunek wykonany w~MATLABie i~wygenerowany skryptem \texttt{matlab2tikz}: prawidłowy dobór kolorów}
\label{r_wykres_kolory1}
\end{figure}

\begin{figure}[H]
\centering
	\input{./rysunki/matlab2tikz/wykres_kolory2.tex}
\caption{Przykładowy rysunek wykonany w~MATLABie i~wygenerowany skryptem \texttt{matlab2tikz}: nieprawidłowy dobór kolorów}
\label{r_wykres_kolory2}
\end{figure}

\section{Lokalizacja rysunków i~tabel}
Zgodnie z~filozofia systemu \LaTeX, wszystkie rysunki i~tabele powinniśmy umieszczać jako ,,pływające", a~więc w otoczeniu \verb+figure+ oraz \verb+table+, możliwe pozycje: \verb+[t]+ -- na górze strony, \verb+[b]+ -- na dole strony, \verb+[p]+ -- na stronie rysunków/tabel (możliwe są również kombinacje tych opcji). Można również zastosować opcję \verb+[h]+ (jeżeli się da, to w~tym miejscu). Wymienione opcje działają dobrze, pod warunkiem jednak, że w~dokumencie znajduje się ,,odpowiednia'' ilość tekstu, wystarczająca do ładnego ułożenia razem z~rysunkami i~tabelami. Niestety, jeżeli rysunków i~tabel jest znacznie więcej niż tekstu, mogą się znajdować daleko od miejsca, w~którym są omawiane. Należy wówczas zastosować opcję \verb+[H]+ z~pakietu \verb+float+ (dokładnie w~tym miejscu). Wadą takiego rozwiązania jest prawdopodobieństwo wystąpienia pustych fragmentów stron, ale rysunki i~tabele są dokładnie w~tym miejscu, w~którym wskażemy. Aby uniknąć umieszczenia rysunków w~tekście kolejnego rozdziału można również zastosować polecenie \verb+\FloatBarrier+ z~pakietu \verb+placeins+.


